\section{Conclusion}\label{sec:conclusion}

This study measured how the settlement error of demand-response customer baselines depends on the hours each baseline reads, using {{sample.n_buildings|int}} buildings at {{sample.n_sites|int}} campus-heavy sites in the 2017 held-out year, hot weekdays as proxy events and injected load changes. All five confirmatory decisions met their Holm-adjusted criteria, with H3 tested for equivalence.

RQ1. Day-of adjustment reduced the hot-day error of PJM, and the hot-day bias of unadjusted PJM was more negative than its random-weekday bias; both are replications. Under the signed-rank location (symmetry) assumption, CG-adj was equivalent to PJM-adj at the site median within the pre-set margin, but {{derived.H3.n_sites_outside_margin|int}} of {{sample.n_sites|int}} sites lay outside it and the conditional bootstrap interval crossed the lower margin.

RQ2. In the four-hour injected pre-event stress test, PJM-adj lost its accuracy advantage over SCM-ridge at a break-even pre-event increase of {{H4.median_site_pcross|pct1}} at the median site, the median of the site break-evens with a site-level sign test. This is a worst-case placement bound, because the shaped hours cover all three adjustment hours. SCM-ridge is immune by construction to the specified four-hour injection. The two-hour result is building-pooled, with no site-level test, and its interval ({{pcross.pjmA.2h:ci_low|pct1}} to {{pcross.pjmA.2h:ci_high|pct1}}) includes the pre-set ten percent threshold.

RQ3. In the injected repeated-dispatch stress test, PJM-adj had the lower absolute error at each tested dispatch count in both the 14:00 to 18:00 and 16:00 to 20:00 windows. In the confirmatory 14:00 to 18:00 window, SCM-ridge's absolute error grew faster (slope difference {{H5.eta_diff|.4f}}, interval {{H5.eta_diff:ci_low|.4f}} to {{H5.eta_diff:ci_high|.4f}}), although PJM-adj's also grew. In the pre-specified secondary 16:00 to 20:00 window this slope contrast was inconclusive.

Operators should therefore choose and monitor a baseline by the hours it reads and by the expected load-shaping and dispatch pattern, tracking the break-even pre-event increase and the absolute-error slope against dispatch count. These results come from stress tests on injected load, not from observed participant behavior. The pipeline is open, and its analysis plan is pre-specified and hash-frozen, so other researchers can rerun the stress tests on their own meters and rules.
